\documentclass[10pt,letterpaper]{article}
\usepackage{cornell-notes}

\title{Chapter 6: Filters}
\author{}
\date{}
\setDocTitle{Chapter 6: Filters}
\setDocSubtitle{Electronics}
\setDocOwner{Electronics Cornell Notes}
\setDocDate{\today}
\setCornellCollection{Electronics Cornell Notes}
\setCornellUnitType{Chapter}
\setCornellUnitNumber{6}
\setCornellUnitTitle{Filters}
\hypersetup{pdftitle={Chapter 6: Filters},pdfsubject={Cornell notes},pdfauthor={}}
\begin{document}
\maketitle

\begin{CornellOverviewBox}{Concept}
Filters shape spectra to reject interference, limit noise, separate channels, prevent aliasing, or control waveform response. A usable design begins with passband, stopband, transient, impedance, dynamic-range, and tolerance requirements; selects an approximation and order; then realizes it with passive, active, switched-capacitor, or digital structures whose nonideal behavior is verified.
\end{CornellOverviewBox}
\begin{CornellOverviewBox}{Concept}
\textbf{Prerequisites.} Complex impedance, Bode magnitude and phase, RC/LC time response, op-amp feedback, and logarithmic frequency scales.

\textbf{After studying this chapter, you should be able to:}
\begin{itemize}
  \item Translate system requirements into poles, zeros, order, cutoff, ripple, attenuation, and transient constraints.
  \item Analyze passive RC and LC networks including source/load interaction.
  \item Choose Butterworth, Bessel, Chebyshev, notch, allpass, state-variable, switched-capacitor, or digital approaches.
  \item Evaluate op-amp bandwidth, dynamic range, tolerance sensitivity, noise, and stability in active filters.
\end{itemize}\textbf{Central design questions.}
\begin{itemize}
  \item Is the primary requirement frequency selectivity, waveform fidelity, delay, noise bandwidth, or anti-aliasing?
  \item How will source/load impedance and component spread move poles and $Q$?
  \item Does the chosen amplifier preserve the intended response at signal amplitude and frequency?
\end{itemize}\end{CornellOverviewBox}
\section{Cornell Notes}
\begin{CornellNotesTable}
\CornellNoteRow{6.1: Specification first}{State filter function, passband edge, allowed ripple, stopband edge/attenuation, source and load, signal range, noise, phase or group-delay needs, transient overshoot, power, and component tolerance. Order controls asymptotic selectivity, but approximation controls the trade among flatness, transition width, phase, and ringing. Define whether cutoff means $-3\,\mathrm{dB}$, a passband edge, or another contractual point.}
\CornellNoteRow{6.2.1: RC response}{A first-order low-pass has $H(s)=1/(1+sRC)$; a high-pass has $H(s)=sRC/(1+sRC)$. At $f_c=1/(2\pi RC)$ magnitude is $1/\sqrt{2}$ and phase is $\mp45^{\circ}$. Cascaded passive sections interact unless buffered because each section loads the previous one. The source resistance and load must be part of the transfer function.}
\CornellNoteRow{6.2.2--6.2.3: LC networks}{LC filters can achieve low loss and high-frequency operation, but terminations are essential to the intended response. Inductor series resistance, capacitor ESR/ESL, self-resonance, saturation, and layout alter poles and damping. Energy exchange causes ringing when $Q$ is high. Simulate with parasitics and verify with the actual source/load impedance.}
\CornellNoteRow{6.2.4--6.2.7: Types and realization}{Butterworth emphasizes flat magnitude, Bessel smooth delay and transient response, and Chebyshev sharper transition with ripple. Elliptic responses add finite-frequency zeros and high sensitivity. Normalize a low-pass prototype, scale frequency and impedance, split high order into stable sections, and order stages to protect dynamic range. Active realizations avoid inductors but add amplifier limits.}
\CornellNoteRow{6.3.1--6.3.2: VCVS sections}{A voltage-controlled voltage-source (Sallen--Key) section realizes second-order poles economically. Component ratios and amplifier gain set natural frequency and $Q$. High-$Q$ sections are sensitive to tolerances and finite op-amp gain. Use design tables or verified equations, then Monte Carlo the actual preferred values. The op-amp's noise gain and GBW must substantially exceed the section demands.}
\CornellNoteRow{6.3.3: State-variable filters}{State-variable structures generate simultaneous low-pass, band-pass, and high-pass outputs and tune frequency and $Q$ with relatively independent parameters. They use more amplifiers but can be easier to tune and less sensitive at high $Q$. Integrator offsets can accumulate, so provide dc control and verify output swing in every internal state.}
\CornellNoteRow{6.3.4--6.3.5: Notch and allpass}{Twin-T networks create a notch by cancellation; depth depends strongly on matching while notch frequency depends on products. An active loop can sharpen the notch but may become unstable. Allpass filters preserve ideal magnitude while changing phase/group delay. Use them for phase equalization or quadrature generation, not as a substitute for bandwidth limiting.}
\CornellNoteRow{6.3.6--6.3.8: Switched-capacitor and DSP}{Switched-capacitor filters set ratios using capacitor ratios and clock frequency, offering accurate tuning but adding clock feedthrough, alias zones, dynamic-range limits, and output steps. Digital filters provide reproducibility and programmable response after sampling; they require anti-alias filtering, adequate sample rate, numeric headroom, and reconstruction filtering. Partition analog and digital filtering from end-to-end noise and latency needs.}
\end{CornellNotesTable}
\begin{CornellSummaryBox}{Summary}
Filter design is an exercise in controlled tradeoffs, not merely selecting a cutoff. Frequency magnitude, phase, group delay, transient response, noise bandwidth, impedance, dynamic range, component sensitivity, and implementation limits are coupled. Scale a suitable response, realize it in manageable sections, and verify the real component, amplifier, sampling, and layout effects.
\end{CornellSummaryBox}
\section{Worked Examples}
\begin{CornellExamBox}{First-order anti-alias section}
For $R=7.5\,\mathrm{k\Omega}$ and $C=4.7\,\mathrm{nF}$, $f_c=1/(2\pi RC)=4.52\,\mathrm{kHz}$. At $20\,\mathrm{kHz}$ the ideal attenuation is $20\log_{10}[1/\sqrt{1+(20/4.52)^2}]=-13.1\,\mathrm{dB}$. This is rarely enough as an anti-alias filter by itself; compare attenuation at every alias-producing frequency.
\end{CornellExamBox}
\begin{CornellExamBox}{Second-order damping}
For a second-order denominator $s^2+2\zeta\omega_0s+\omega_0^2$, let $f_0=2\,\mathrm{kHz}$ and $Q=0.707$. Then $\zeta=1/(2Q)=0.707$. The section is near a Butterworth pair: flat magnitude but some step overshoot. A Bessel-oriented section would trade transition sharpness for smoother time response.
\end{CornellExamBox}
\begin{CornellExamBox}{Tolerance shift}
An RC cutoff uses $R$ at $\pm1\%$ and $C$ at $\pm5\%$. Worst-case fractional cutoff shift is approximately $\pm(1\%+5\%)=\pm6\%$ for small errors; RSS is $5.10\%$ only if statistical independence is justified. High-$Q$ response can be far more sensitive than this simple product.
\end{CornellExamBox}
\section{Design and Troubleshooting}
\begin{CornellExamBox}{Design and Troubleshooting}
\begin{enumerate}
  \item Write passband, stopband, phase/delay, transient, noise, impedance, dynamic-range, and tolerance requirements.
  \item Choose approximation and minimum practical order; factor into first/second-order sections and scale frequency/impedance.
  \item Select topology and parts; include source/load, amplifier GBW/noise/slew, parasitics, and preferred-value tolerances.
  \item Simulate nominal, corners, and Monte Carlo; prototype and measure magnitude, phase, step response, noise, and internal clipping.
\end{enumerate}\end{CornellExamBox}
\begin{CornellExamBox}{Symptoms, likely causes, and checks}
\begin{itemize}
  \item Cutoff shifted: source/load resistance, component tolerance, probe loading, or wrong scaling.
  \item Unexpected peaking: excessive $Q$, op-amp phase error, topology sensitivity, or insufficient damping.
  \item Notch too shallow: component ratio mismatch, source/load loading, leakage, or layout coupling around the notch path.
  \item Output clips though input is small: internal section gain and ordering create larger intermediate signals.
  \item Clock tones or aliases: switched-capacitor/DSP clock feedthrough, inadequate analog filtering, or sampling-plan error.
\end{itemize}\end{CornellExamBox}
\begin{CornellWarningBox}{Safety and Limitations}
Passive filters connected to mains, power converters, or high-energy sources need voltage-rated capacitors, current-rated inductors, discharge paths, creepage, and flame-safe components. Inductors can saturate and overheat; capacitors can retain charge. Verify measurement equipment grounding before probing line-referenced networks.
\end{CornellWarningBox}
\section{Key-Equation Sheet}
\begin{CornellExamBox}{Key Equations}
\begin{itemize}
  \item $f_c=1/(2\pi RC)$ for a first-order RC pole or zero.
  \item $H_{LP}(s)=1/(1+s/\omega_c)$ and $H_{HP}(s)=(s/\omega_c)/(1+s/\omega_c)$.
  \item Second-order denominator: $s^2+(\omega_0/Q)s+\omega_0^2$; $\zeta=1/(2Q)$.
  \item Ideal resonant frequency: $f_0=1/(2\pi\sqrt{LC})$.
  \item A cascade multiplies transfer functions and adds dB magnitudes: $H_{tot}=\prod H_k$.
  \item For white input noise density, equivalent noise bandwidth -- not only the $-3\,\mathrm{dB}$ point -- sets integrated noise.
\end{itemize}\end{CornellExamBox}
\section{Study and Review}
\subsection{Design checklist}
\begin{itemize}
  \item Cutoff, ripple, stopband, phase/delay, and transient meanings are explicit.
  \item Source/load impedances are included in synthesis and test.
  \item Section $Q$, tolerance sensitivity, dynamic range, and order are acceptable.
  \item Amplifier GBW, slew, noise, output drive, and stability are adequate.
  \item Analog anti-alias/reconstruction needs are included around sampled filtering.
\end{itemize}\subsection{Common misconceptions}
\begin{itemize}
  \item A cutoff frequency alone does not specify a filter.
  \item Cascaded passive RC sections do not multiply ideally when they load each other.
  \item Higher order is not automatically better; it adds delay, sensitivity, cost, and ringing.
  \item A maximally flat magnitude response is not maximally flat delay.
  \item Digital filtering cannot remove aliases created before sampling.
\end{itemize}\subsection{Active-recall questions}
\begin{enumerate}
  \item What distinguishes Butterworth, Bessel, and Chebyshev responses?
  \item Why does termination matter to an LC filter?
  \item How are $Q$ and damping ratio related?
  \item Why can section order affect clipping?
  \item What determines notch depth?
  \item Which op-amp parameter most directly erodes high-$Q$ accuracy?
  \item Why use a state-variable structure?
  \item What artifacts accompany switched-capacitor filters?
  \item What is equivalent noise bandwidth?
  \item Why is an analog anti-alias filter still required?
\end{enumerate}\subsection{Original practice problems}
\begin{enumerate}
  \item Find $C$ for a $1.5\,\mathrm{kHz}$ RC low-pass with $R=12\,\mathrm{k\Omega}$.
  \item A two-pole low-pass is $-40\,\mathrm{dB/decade}$ asymptotically. Estimate change from $10\,\mathrm{kHz}$ to $100\,\mathrm{kHz}$ far above its poles.
  \item For $L=2.2\,\mathrm{mH}$ and $C=47\,\mathrm{nF}$, estimate ideal resonance.
\end{enumerate}\subsection{Answer key / solution outline}
\begin{enumerate}
  \item $C=1/(2\pi Rf_c)=8.84\,\mathrm{nF}$; choose a realizable value and recompute tolerance range.
  \item Approximately $-40\,\mathrm{dB}$ over one decade.
  \item $f_0=1/(2\pi\sqrt{LC})=15.7\,\mathrm{kHz}$; parasitics and loading set actual response.
\end{enumerate}\begin{CornellSummaryBox}{Summary}
Specify the complete response, choose an approximation deliberately, realize it in well-behaved sections, and verify loading, component spread, amplifier limits, internal dynamic range, sampling interactions, and transient response.
\end{CornellSummaryBox}

\end{document}
